\begin{table}[t]
\centering\small
\begin{tabular}{@{}lrrr@{}}
\toprule
outcome method & $n$ & complete & partial \\
\midrule
key from ciphertext only & 50 & 12 & 33 \\
key from external plaintext & 18 & 5 & 13 \\
key from adjacent plaintext & 68 & 22 & 43 \\
\quad\emph{key recovered, all routes} & \emph{136} & \emph{39} & \emph{89} \\
external key matched & 38 & 11 & 21 \\
known key applied & 95 & 43 & 42 \\
existing decipherment & 109 & 91 & 17 \\
\midrule
not solved & 122 & -- & -- \\
\midrule
not applicable & 50 & -- & -- \\
\midrule
all targets & 550 & 184 & 169 \\
\bottomrule
\end{tabular}
\caption{\label{tab:methods} Outcome method and extent for the 550 targets. The first three methods recover the key here; the next two apply a key that already existed; the sixth transcribes and checks a decipherment that already existed. Complete means the read bar of Section~\ref{sec:schema} is met.}
\end{table}

\begin{table}[t]
\centering\small
\begin{tabular}{@{}lr@{\qquad}lr@{}}
\toprule
century & $n$ & plaintext language & $n$ \\
\midrule
14th c. & 5 & French & 133 \\
15th c. & 37 & English & 99 \\
16th c. & 190 & Italian & 90 \\
17th c. & 108 & German & 59 \\
18th c. & 110 & Spanish & 50 \\
19th c. & 41 & Latin & 39 \\
20th c. & 32 & Dutch & 15 \\
undated & 27 & Polish & 11 \\
\bottomrule
\end{tabular}
\caption{\label{tab:corpus} The corpus by century of the earliest document and by plaintext language (a target with letters in two languages counts under both; 20 languages in all).}
\end{table}

\begin{table}[t]
\centering\small
\resizebox{\columnwidth}{!}{%
\begin{tabular}{@{}lrrrrr@{}}
\toprule
starting point & $n$ & recovered & applied & existed & not solved \\
\midrule
ciphertext only & 125 & 46 & 6 & 4 & 69 \\
crib from clear text & 35 & 26 & 2 & 3 & 4 \\
decipherment or known plaintext & 115 & 48 & 8 & 50 & 9 \\
a key (archive, print, sibling) & 151 & 14 & 116 & 13 & 8 \\
none recorded & 74 & 2 & 1 & 39 & 32 \\
\midrule
all & 500 & 136 & 133 & 109 & 122 \\
\bottomrule
\end{tabular}}
\caption{\label{tab:attack} Outcome by what the attack started from (profile field conditions.attack). Columns: key recovered (from ciphertext only, adjacent or external plaintext); key applied (a known key, or an external key matched to the document); a decipherment already existed; not solved. Targets classed not applicable are left out.}
\end{table}

\begin{table}[t]
\centering\small
\resizebox{\columnwidth}{!}{%
\begin{tabular}{@{}lrrrrr@{}}
\toprule
prior reading & $n$ & recovered & applied & existed & not solved \\
\midrule
none known & 152 & 56 & 28 & 0 & 68 \\
existed, known before the attempt & 183 & 42 & 57 & 68 & 16 \\
existed, found during the attempt & 103 & 21 & 34 & 39 & 9 \\
existed, found after or not at all & 25 & 12 & 9 & 1 & 3 \\
existed, when found unknown & 2 & 0 & 1 & 1 & 0 \\
unknown & 35 & 5 & 4 & 0 & 26 \\
\midrule
all & 500 & 136 & 133 & 109 & 122 \\
\bottomrule
\end{tabular}}
\caption{\label{tab:prior} Outcome by whether a reading of the text existed anywhere before this project, and when it was found (the contamination question). Columns as in Table~\ref{tab:attack}.}
\end{table}
